Let \(s \in S\), and let \(f, g, h \in M(S,G)\).
Since \((f+g)(x)\) has the unique value \(f(x)+g(x) \in G\),
the sum \(f+g\) is a well-defined function in \(M(S,G)\).

Moreover,
\[
((f+g)+h)(x)
= (f+g)(x)+h(x)
= (f(x)+g(x))+h(x)
= f(x)+(g(x)+h(x))
= f(x)+(g+h)(x)
= (f+(g+h))(x),
\]
and hence the operation is associative.
Similarly, if \(G\) is abelian, then \(f+g = g+f\).

Let \(0_S \in M(S,G)\) be the function defined by \(x \mapsto e\),
and let \(-f \in M(S,G)\) be the function defined by
\(x \mapsto f(x)^{-1}\).
It is clear that \(0_S\) is an identity element and that \(-f\) is an inverse
of \(f \in M(S,G)\).
